\chapter{Aspectos Conceituais} \label{cap:conceitos}

Este capítulo reúne os conceitos sobre os quais o restante do documento se
apoia e, junto com eles, a revisão da literatura. Os dois conteúdos aparecem
entrelaçados deliberadamente, uma vez que quase toda definição empregada aqui
provém de um trabalho específico e carrega as escolhas que esse trabalho fez. A
Seção~\ref{sec:interrupcoes} descreve o modelo de execução de programas
orientados a interrupções. A Seção~\ref{sec:defeitos} define o defeito que se
busca encontrar. A Seção~\ref{sec:tecnicas} apresenta as técnicas de análise
empregadas pela literatura e por este trabalho, e a Seção~\ref{sec:ferramentas}
descreve as quatro ferramentas que formam o estado da arte. Os
\textit{benchmarks} vêm na Seção~\ref{sec:benchmarks-conceito}, as métricas de
avaliação na Seção~\ref{sec:soundness} e o uso de modelos de linguagem em
análise estática na Seção~\ref{sec:llm}.

Uma advertência vale para todo o capítulo. Salvo indicação em contrário, os
números citados são os relatados pelos próprios autores. Muitos foram medidos
em hardware diferente, sobre contagens diferentes do mesmo \textit{benchmark},
ou copiados de outros artigos porque as ferramentas de comparação não estavam
disponíveis. Eles descrevem o que cada trabalho afirma, não um resultado
comparável entre trabalhos.

\section{Programas orientados a interrupções}
\label{sec:interrupcoes}

\subsection{Fluxos, prioridades e preempção assimétrica}

Um programa orientado a interrupções é formado por um fluxo principal, que
executa continuamente, e por um conjunto de rotinas de tratamento de
interrupção. A literatura o descreve como
$P = \mathit{Main} \parallel \mathit{ISR}_1 \parallel \dots \parallel
\mathit{ISR}_n$, com uma função que atribui uma prioridade a cada fluxo
\cite{bmc4av}. Cada fluxo é uma raiz de execução, no sentido de que tudo o que
ele alcança por chamadas de função pode executar sob sua identidade e com sua
prioridade.

A preempção entre esses fluxos é assimétrica. Uma \isr\ interrompe o fluxo
principal, mas o fluxo principal nunca interrompe uma \isr; uma \isr\ de
prioridade alta interrompe uma de prioridade baixa, e nunca o contrário. Entre
\textit{threads}, a preempção é simétrica, e é dessa diferença que
\citeonline{sdracer} tiram as três consequências que separam os dois mundos.
Uma \isr\ não pode bloquear, já que executa até o fim a menos que outra de
prioridade maior a interrompa. A relação \textit{happens-before}, base dos
detectores de corrida entre \textit{threads}, deixa de ser correta quando uma
interrupção ocorre durante o próprio cálculo dessa relação. E a sincronização
não é feita com travas, mas com mascaramento de interrupções. As técnicas de
\textit{threads}, portanto, não se transferem, e a literatura de interrupções
raciocina sobre padrões de acesso e prioridades.

A convenção de numeração das prioridades diverge entre os artigos, o que exige
atenção de quem reimplementa qualquer um deles.
\citeonline{sdracer} e \citeonline{intrace} formalizam números maiores como
prioridades menores, e \citeonline{bmc4av} adotam o contrário. O
\textit{benchmark} usado por todos eles declara explicitamente que número maior
significa prioridade maior \cite{racebench}, e é essa a convenção adotada aqui.

\subsection{Aninhamento}

Uma \isr\ interrompida por outra de prioridade maior, que por sua vez pode ser
interrompida, forma uma cadeia de aninhamento (\textit{nesting}). O aninhamento
multiplica as intercalações possíveis e é a característica que mais distingue
as ferramentas entre si. O NIChecker o trata nativamente, codificando
prioridade e mascaramento nas condições de invocação de cada \isr\
\cite{nichecker}. O IntRace o aborda por uma conversão sequencial de dentro
para fora \cite{intrace}, o BMC4AV por relações de ordem parcial entre acessos
\cite{bmc4av}, e o SDRacer exclui interrupções reentrantes do seu
modelo \cite{sdracer}.

\subsection{Sincronização por mascaramento}

Sem travas, o programador impede que a interrupção ocorra. Ele mascara a
interrupção antes do trecho que precisa ser indivisível e a reabilita depois.
No Racebench isso é feito com \texttt{disable\_isr(n)} e
\texttt{enable\_isr(n)}, em que \texttt{n = -1} designa todas as interrupções
\cite{racebench}. Em código real, o mascaramento pode vir de APIs padronizadas,
como \texttt{disable\_irq} no Linux, de escritas diretas em registradores de
controle ou de variáveis mantidas pelo próprio programa. \citeonline{intrace}
resolvem o problema pedindo ao usuário um arquivo de configuração que enumere
os mecanismos não padronizados da plataforma.

Dois cuidados decidem se uma análise de mascaramento perde defeitos. O primeiro
é que uma interrupção só pode ser dada como mascarada num ponto se estiver
mascarada em todos os caminhos que chegam àquele ponto, pois ``mascarada em
algum caminho'' não oferece proteção alguma. O segundo é menos evidente e aparece no
próprio \textit{benchmark}. O mascaramento feito por um fluxo pode ser desfeito
por outro. No caso \texttt{svp\_simple\_001\_001} do Racebench, o fluxo
principal mantém a interrupção 2 mascarada durante todo o intervalo de um
defeito anotado, e o defeito ainda assim é real, porque a interrupção 1
continua habilitada, interrompe esse intervalo e reabilita a interrupção 2 por
dentro. Uma análise que considerasse apenas o fluxo local
descartaria um defeito verdadeiro sem qualquer indício de erro.

\section{O defeito de concorrência tratado neste trabalho}
\label{sec:defeitos}

A literatura enuncia o problema de duas maneiras. SDRacer e IntRace falam em
\textit{data race} e o definem sobre dois acessos, um de cada fluxo, sendo ao
menos um deles uma escrita, perguntando se a \isr\ pode executar no instante
exato do acesso do outro fluxo \cite{sdracer,intrace}. NIChecker e BMC4AV falam
em violação de atomicidade e a definem sobre três acessos, dois do fluxo
interrompido e um da \isr\ entre eles, perguntando se a execução resultante
equivale a alguma execução em série dos fluxos \cite{nichecker,bmc4av}. Este
trabalho adota a segunda formulação, por duas razões. O gabarito do Racebench
anota todos os seus defeitos sobre três acessos, inclusive nos casos usados
pelos trabalhos da primeira linha, e a pergunta sobre um intervalo é
estritamente mais exigente que a pergunta sobre um instante. Os dois termos são
empregados como sinônimos no restante do documento.

Uma violação de atomicidade envolve três acessos à mesma posição de memória.
Dois deles, $A_1$ e $A_2$, pertencem ao fluxo interrompido e o programador os
supunha indivisíveis. O terceiro, $B$, pertence a uma \isr\ e ocorre entre os
outros dois. Das oito combinações possíveis de leitura (R) e escrita (W) sobre
a tripla, quatro produzem um estado que nenhuma execução em série produziria.
O Quadro~\ref{quadro:padroes} mostra cada uma delas.

\begin{quadro}[htb]
  \caption{\label{quadro:padroes}Os quatro padrões não serializáveis}
  \begin{tabular}{@{}lL{9.5cm}@{}}
    \toprule
    Padrão $(A_1, B, A_2)$ & Por que quebra a atomicidade \\
    \midrule
    $(R,W,R)$ & as duas leituras do mesmo fluxo veem valores diferentes \\
    $(R,W,W)$ & $A_2$ escreve com base numa leitura que ficou obsoleta \\
    $(W,R,W)$ & a \isr\ enxerga um valor intermediário que não deveria existir \\
    $(W,W,R)$ & $A_2$ lê o valor escrito pela \isr, e não o escrito por $A_1$ \\
    \bottomrule
  \end{tabular}
  \fonte{Elaborado pelos autores a partir de \citeonline{nichecker} e
  \citeonline{bmc4av}.}
\end{quadro}

As outras quatro combinações, $(R,R,\ast)$, $(W,R,R)$ e $(W,W,W)$, são
serializáveis. O padrão $(R,W,W)$ é disputado, e a divergência importa para
qualquer contagem. \citeonline{bmc4av} o classificam como benigno e excluem do
Racebench os casos que só o contêm, enquanto \citeonline{nichecker} o contam
como defeito. Este trabalho o mantém, porque o gabarito do Racebench o anota em
10 dos 48 defeitos que contamos nos casos simples.

Nem todo defeito tem consequência, e a literatura separa as duas perguntas.
\citeonline{intrace} adotam o critério de que um defeito é nocivo quando a
variável compartilhada alimenta uma condição de desvio ou serve de índice de
vetor ou de ponteiro. Os mesmos autores relatam corridas deixadas sem proteção
deliberadamente sem proteção, por razões de tempo real, como uma de suas duas
fontes conhecidas de falsos positivos. Julgar a nocividade é julgar a intenção do programa, o que
é distinto de decidir se a intercalação é possível, e as duas perguntas são
mantidas separadas ao longo de todo este trabalho.

\section{Técnicas de análise de programas}
\label{sec:tecnicas}

\subsection{Análise estática}

A análise estática raciocina sobre o código sem executá-lo, e três estruturas
sustentam o que se faz aqui. O grafo de chamadas indica quais funções cada
função pode chamar e, por extensão, tudo o que é alcançável a partir de cada
fluxo. O grafo de fluxo de controle (\textit{Control-Flow Graph}, CFG) indica,
dentro de uma função, em que ordem as instruções podem executar. A análise de
ponteiros indica a que posições de memória cada ponteiro pode apontar. Ela é do
tipo \textit{may} quando superaproxima as possibilidades e do tipo \textit{must}
quando as subaproxima. Para não perder defeitos é preciso usar a forma
\textit{may}, e uma posição é considerada compartilhada quando os conjuntos de
posições que dois fluxos podem alcançar se interceptam.

Este trabalho faz a análise sobre a representação intermediária do compilador
LLVM \cite{llvm} e usa a biblioteca SVF na análise de ponteiros \cite{svf}. Os
metadados de depuração do LLVM, em particular o \texttt{DILocation}, permitem
associar cada instrução analisada a uma linha do código-fonte C. Sem esse
vínculo, nenhum revisor, humano ou modelo de linguagem, conseguiria interpretar
o resultado.

\subsection{Execução simbólica e viabilidade de caminho}

Na execução simbólica as entradas do programa são tratadas como símbolos, e
cada caminho percorrido acumula uma condição sobre esses símbolos. Um
solucionador decide se a condição é satisfazível, isto é, se existe alguma
entrada que leve o programa por aquele caminho. \citeonline{sdracer} usam a
técnica para gerar entradas e intercalações que alcancem os pontos suspeitos, e
registram uma limitação relevante, a de que registradores de hardware não
constituem variáveis livres, pois seus valores dependem do modelo do
dispositivo.

A análise de viabilidade de caminho aplica a mesma ideia a um candidato
específico. \citeonline{intrace} constroem um resumo simbólico da \isr,
inserem-no no ponto de preempção e entregam as restrições resultantes ao
solucionador Z3 \cite{z3}. Quando o solucionador responde que as restrições são
insatisfazíveis, a intercalação é impossível, e esse veredicto é uma prova.
Essa propriedade ocupa um lugar central no projeto apresentado adiante, porque
só uma prova de impossibilidade autoriza descartar um candidato.

\subsection{Verificação de modelos limitada}

A verificação de modelos limitada (\textit{Bounded Model Checking}, BMC)
procura defeitos em todas as execuções de comprimento até um limite $u$,
codificando o programa como uma fórmula para um solucionador SAT ou SMT. Sua
vantagem é produzir um contraexemplo, isto é, uma execução concreta que exibe o
defeito. Sua limitação é o crescimento exponencial do espaço de estados com
$u$, e alguns casos do Racebench exigem desenrolar um laço dez mil vezes antes
de o defeito aparecer \cite{nichecker}. A ferramenta de referência dessa
família é o CBMC \cite{cbmc}.

Duas técnicas atenuam o problema. A sequencialização preguiçosa (\textit{lazy
sequentialization}) traduz o programa concorrente num programa sequencial não
determinístico que simula suas execuções e o entrega a um verificador
sequencial \cite{lazycseq}. A abstração de laços substitui laços por versões
superaproximadas antes do desenrolamento, para que defeitos profundos apareçam
com um limite pequeno \cite{nichecker}. Em qualquer caso, a ausência de
contraexemplo vale apenas até o limite escolhido e não é prova de ausência de
defeitos.

\section{Ferramentas de detecção para programas com interrupções}
\label{sec:ferramentas}

\subsection{SDRacer}

\citeonline{sdracer} combinam análise estática, execução simbólica e simulação
numa plataforma virtual para detectar, validar e reparar condições de corrida
em software embarcado. A ferramenta trabalha em quatro etapas. A análise
estática emite os alertas iniciais. A execução simbólica gera entradas e
intercalações que alcancem cada alerta e elimina os de condições de caminho
conflitantes. A plataforma virtual Simics força as interrupções a ocorrer nos
pontos suspeitos e mantém apenas os candidatos que de fato se intercalam. Por fim, a
ferramenta sugere correções, e o vocabulário dessas correções é específico de
interrupção, incluindo desabilitar e reabilitar a interrupção em volta do
acesso, inserir travas, estender uma seção crítica existente ou fundir seções
adjacentes.

Os autores relatam 190 condições de corrida detectadas em nove
\textit{benchmarks} embarcados, com a execução simbólica eliminando 40,3\% dos
alertas estáticos e a validação dinâmica mais 36,7\%. As correções custaram
menos de 0,09 de sobrecarga em 9 dos 11 programas avaliados; nos outros dois a
sobrecarga foi bem maior, porque desabilitar interrupções alterou o fluxo de
controle da tarefa principal. A ausência de falsos negativos foi afirmada por
inspeção manual. É a única ferramenta revisada que repara defeitos, e a
contrapartida é que sua detecção se limita ao que a plataforma virtual consegue
simular.

\subsection{IntRace}

\citeonline{intrace} detectam corridas de forma puramente estática, em três
estágios de custo e precisão crescentes. O primeiro constrói o CFG a partir do
Clang e de uma passagem LLVM, identifica os recursos compartilhados por análise
de ponteiros interprocedural e casa acessos, ignorando deliberadamente
condições de desvio e estado das interrupções para não perder nada. O segundo
descarta candidatos cujos fluxos não podem se sobrepor, seja por prioridade,
seja por mascaramento. O terceiro verifica a viabilidade de cada candidato
restante com o Z3.

A medição por estágio é o resultado mais aproveitável do artigo. São cerca de
208 candidatos por programa real após o primeiro estágio; o segundo elimina
54,2\% deles; o terceiro elimina mais 86,2\% do que lhe chega e consome 69,7\%
do tempo total de análise. Sobre 30 dos 31 casos simples do Racebench, os
autores relatam 5 falsos positivos e 73,2\% a menos de falsos positivos que a
ferramenta Rchecker, embora os números do Rchecker tenham sido copiados do
artigo original, pois a ferramenta não estava disponível. Sobre nove programas
industriais, relatam 118 corridas e 9 falsos positivos. A arquitetura em três
estágios do IntRace é o ponto de partida direto da etapa estática deste
trabalho.

\subsection{NIChecker}

\citeonline{nichecker} verificam violações de atomicidade traduzindo o programa
com interrupções num programa sequencial, por sequencialização preguiçosa, e
entregando o resultado ao CBMC \cite{lazycseq,cbmc}. As condições de invocação
de cada \isr\ no programa sequencial codificam prioridade e mascaramento, o que
dá suporte nativo ao aninhamento. Antes do desenrolamento a ferramenta aplica
abstração de laços, o que lhe permite encontrar defeitos que exigiriam
desenrolar um laço dez mil vezes, e o fatiamento de programa somado à redução
de pontos de preempção acelera a verificação em até 42,2\%. É o único trabalho
revisado que prova a correção limitada de sua tradução.

Os autores relatam 96,4\% de precisão no Racebench e 47 violações encontradas,
sem falsos positivos, numa suíte de 18 programas reais. A limitação prática mais
séria é operacional, já que o usuário precisa indicar, a cada execução, qual
variável global e qual padrão verificar.

\subsection{BMC4AV}

\citeonline{bmc4av} também partem do CBMC, mas, em vez de sequencializar o
programa, guiam o solucionador com um grafo de acessos à memória construído
durante a resolução. Primeiro identificam violações potenciais por casamento de
padrões e, para cada uma, as relações de ordem que precisariam valer; depois
estendem o grafo guiados por essas relações até confirmar ou descartar a
violação. É a única das quatro ferramentas revisadas cujo código-fonte está
disponível.

Sua contribuição mais relevante para este trabalho não é o resultado que
reporta, e sim a crítica que dirige ao NIChecker, segundo a qual uma ferramenta
com precisão de 100\% encontraria menos de 40\% das violações existentes na
suíte de 18 programas reais. A crítica é aqui tomada como sintoma do problema
de avaliação discutido na Seção~\ref{sec:soundness}, com duas ressalvas. O
trabalho é um \textit{preprint} sem revisão por pares, e a contagem que sustenta
o número adota uma unidade diferente da do gabarito distribuído com a suíte, o
que explica parte da diferença.

\subsection{Outras ferramentas citadas}

Os trabalhos acima se comparam a ferramentas que não foram revisadas
diretamente aqui. O Rchecker detecta corridas com o CBMC e uma lista de
máscaras de interrupção; o intAtom detecta violações de atomicidade de forma
estática; o CPA4AV usa árvores de alcançabilidade abstratas e falha nos casos
com laços profundos; e o iCBMC estende o CBMC para interrupções aninhadas.
Rchecker e intAtom não puderam ser obtidos pelos autores que os usaram como
comparação, e seus números foram copiados de outros artigos, medidos em outro
hardware \cite{intrace,nichecker,bmc4av}.

\subsection{Comparação}

O Quadro~\ref{quadro:capacidades} resume as quatro ferramentas.

\begin{quadro}[htb]
  \caption{\label{quadro:capacidades}Capacidades das ferramentas revisadas}
  \small
  \begin{tabular}{@{}L{2.4cm}L{2.5cm}L{2.5cm}L{3.0cm}L{2.7cm}@{}}
    \toprule
    & SDRacer & IntRace & NIChecker & BMC4AV \\
    \midrule
    Ano, veículo & 2020, TSE & 2024, J. Supercomp. & 2025, TOSEM & 2026, \textit{preprint} \\
    Método & estático, simbólico e dinâmico & estático e Z3 & sequencialização e BMC & BMC com grafo de acessos \\
    Aninhamento & exclui reentrância & conversão sequencial & nativo & por ordens parciais \\
    Automação & automática & arquivo de configuração & variável e padrão por execução & automática, segundo os autores \\
    Repara defeitos & sim & não & não & não \\
    Código disponível & não & não & não & sim \\
    \bottomrule
  \end{tabular}
  \fonte{Elaborado pelos autores a partir de \citeonline{sdracer},
  \citeonline{intrace}, \citeonline{nichecker} e \citeonline{bmc4av}.}
\end{quadro}

Duas leituras do quadro orientam este trabalho. A primeira é que a ausência de
falsos negativos é afirmada, não medida. SDRacer e IntRace a apoiam em inspeção
manual, e os trabalhos de verificação limitada, na ausência de contraexemplo
até um limite de desenrolamento. Nenhum deles declara as premissas sob as quais
nenhum defeito é perdido, o que torna a afirmação impossível de verificar de
fora. A segunda é que a automação é o eixo real da competição entre eles. Cada
ferramenta nova ataca o trabalho manual exigido pela anterior, e é nesse sentido
que a usabilidade aparece, mais adiante, como parte da contribuição deste
trabalho.

\section{Benchmarks}
\label{sec:benchmarks-conceito}

\subsection{Racebench}

O Racebench \cite{racebench} é o \textit{benchmark} sobre o qual IntRace,
NIChecker e BMC4AV são avaliados. Foi criado para uma competição de protótipos
de detecção de condições de corrida realizada em Hangzhou, na China, em
novembro de 2019, e seus programas são modelados em software embarcado do setor
aeroespacial. A versão 2.1 tem 31 casos simples e 2 complexos, cada um com um
fluxo principal e de uma a três \isr s, em até três níveis de prioridade.

O gabarito é anotado no próprio código, e cada defeito anotado registra os três
acessos que o compõem. Além dos defeitos, o \textit{benchmark} contém armadilhas,
que são falsos positivos inseridos intencionalmente, vários deles trechos
corretamente protegidos por mascaramento. A regra de pontuação da competição
premiava cada detecção com 2 pontos e punia cada falso positivo com 3.

Contando diretamente as anotações dos 31 casos simples, obtivemos 48 defeitos e
38 armadilhas. Os artigos usam outras contagens para os mesmos casos, sendo 50
em \citeonline{intrace}, 54 em \citeonline{nichecker} e 38 sobre um subconjunto
de 25 casos em \citeonline{bmc4av}. Parte da divergência tem explicação
mecânica. As anotações usam ao menos quatro gramáticas incompatíveis e contêm
erros de digitação, de modo que um leitor estrito encontra apenas 28 defeitos
sem acusar nenhum erro. A divergência tem consequência direta sobre qualquer
taxa de recall publicada, uma vez que é o denominador dessa taxa que está em
disputa. Por esse motivo, este trabalho fixa e publica sua unidade de contagem
antes de realizar qualquer medição, conforme descrito na
Seção~\ref{sec:avaliacao}.

\subsection{Suíte de programas reais}

A segunda suíte, usada por NIChecker e BMC4AV, tem 18 programas extraídos de
seis pacotes, entre eles o \textit{firmware} de um registrador de temperatura,
um controle de LEDs, um sistema de freio eletrônico gerado a partir de um
modelo Simulink e três \textit{drivers}, num total de 10.633 linhas
\cite{nichecker,bmc4av}. Cada programa acompanha um arquivo com as prioridades
dos fluxos e um arquivo de gabarito, que registra 45 entradas de violação
verdadeira e 14 de falso positivo.

\section{Correção, recall e precisão}
\label{sec:soundness}

Um falso positivo é um alerta que não corresponde a um defeito real, e um falso
negativo é um defeito real que não gerou alerta. Uma análise é \textit{sound},
ou correta, em relação a um conjunto de premissas quando não produz falsos
negativos sob essas premissas. A ausência incondicional de falsos negativos não
é alcançável para C com ponteiros, tabelas de vetores de interrupção e fluxo de
controle dependente de hardware. O que se pode afirmar honestamente é a
ausência relativa a uma lista explícita de premissas, e é essa a afirmação que
este trabalho se propõe a sustentar.

O Quadro~\ref{quadro:metricas} resume as métricas usadas na literatura e aqui.

\begin{quadro}[htb]
  \caption{\label{quadro:metricas}Métricas de avaliação}
  \small
  \begin{tabular}{@{}L{3.0cm}L{6.4cm}L{4.4cm}@{}}
    \toprule
    Métrica & Definição & O que responde \\
    \midrule
    Precisão & $TP / (TP + FP)$ & quanto do que foi reportado é real \\
    Recall (taxa de acerto) & $TP / (TP + FN)$ & quanto do que existe foi encontrado \\
    \textit{Inspection Ratio} & $(TP + FP) / (TP + FP + TN + FN)$ & quanto é preciso ler \\
    \bottomrule
  \end{tabular}
  \fonte{Elaborado pelos autores a partir de \citeonline{reducing-fp}.}
\end{quadro}

A prioridade dada aqui ao recall tem uma razão concreta. A precisão, sozinha,
pode ser melhorada descartando candidatos, e uma ferramenta pode exibir
precisão perfeita examinando apenas o que o usuário já apontou, perdendo a
maioria dos defeitos. O \textit{Inspection Ratio} não tem essa fragilidade,
porque mede a fração do conjunto que um revisor precisa ler até encontrar todos
os defeitos reais. Com o recall mantido em 100\% por construção, é o
\textit{Inspection Ratio} que mostra o custo real da ferramenta
\cite{reducing-fp}.

\section{Modelos de linguagem aplicados à análise estática}
\label{sec:llm}

\subsection{O que é e o que falha}

Um modelo de linguagem de grande escala (\llm) é uma rede neural treinada para
continuar textos, capaz de ler e produzir código. A instrução dada a ele é
chamada de \textit{prompt}, e para uso programático pede-se uma saída
estruturada, num formato fixo como JSON, que possa ser lida e validada
automaticamente.

Três limitações importam aqui. Um \llm\ pode alucinar, afirmando com confiança
fatos falsos sobre o código. Suas respostas variam entre execuções, de modo que
a mesma pergunta pode receber veredictos diferentes. E seu uso tem custo,
medido em \textit{tokens}, com \citeonline{llift} relatando cerca de 7.000
\textit{tokens} e US\$~0,43 por candidato analisado, com o GPT-4 e os preços de
2024. Qualquer desenho que dependa do julgamento do modelo como decisão final
herda essas três limitações.

\subsection{Triagem de alertas}

\citeonline{llift} partem do UBITect, analisador que procura usos de variáveis
antes da inicialização no \textit{kernel} Linux. A análise estática do UBITect
gera cerca de 140 mil candidatos, e sua etapa de execução simbólica abandona 53
mil deles, ou 40\%, por limite de tempo ou memória. O LLift entrega ao \llm\
apenas esses casos não decididos, com uma política conservadora segundo a qual
tudo o que não for provado seguro é reportado. O trabalho encontrou quatro
defeitos reais confirmados pela comunidade do Linux. Sua ablação mostra o
recall passando de 0,15, com uma pergunta direta, para 1,00 com quatro
componentes de \textit{prompt}, descritos adiante. É a referência estrutural
mais próxima deste trabalho, e seu formato de ablação é adotado aqui.

\citeonline{chatgpt-static} avaliam modelos GPT na remoção de falsos positivos
do analisador Infer e obtêm precisão de 93,88\% para desreferência de ponteiro
nulo, mas apenas 63,33\% para vazamento de recursos, com os mesmos
\textit{prompts} e o mesmo modelo. A lição é que a qualidade da triagem depende
da classe de defeito e não se transfere por suposição, razão pela qual a
aplicação a defeitos de concorrência precisa ser medida, e não presumida.

\citeonline{reducing-fp} filtram alertas do CppCheck com um modelo CodeBERT
ajustado. No conjunto Big-Vul a precisão sobe de 5,88\% para 96,20\%, mas o
recall cai de 46,90\% para 40,66\%, e os próprios autores reconhecem que a
filtragem introduz falsos negativos. É desse trabalho que vem a métrica
\textit{Inspection Ratio}, que no mesmo experimento cai de 50,49\% para 2,44\%.

\subsection{Inferência de especificações}

\citeonline{iris} usam um \llm\ para inferir quais funções de bibliotecas são
fontes e sumidouros de dados contaminados, alimentando com essas
especificações o analisador CodeQL. Num conjunto de 120 vulnerabilidades reais
em Java a detecção sobe de 27, com o CodeQL sozinho, para 55. A precisão
continua baixa, com taxa média de falsas descobertas de 84,82\%, e a triagem
por \llm\ só melhora a precisão com modelos grandes, piorando com modelos
pequenos.

\citeonline{adataint} seguem a mesma ideia, combinando inferência de
especificações com um classificador que filtra alertas. Os resultados vêm de
\textit{benchmarks} sintéticos, e o filtro descarta alertas com recall de
75,4\%, ou seja, cerca de um quarto dos defeitos é perdido por construção.

\subsection{Detecção e reparo}

\citeonline{skipanalyzer} encadeiam três componentes com \llm, fazendo
detecção, filtragem de falsos positivos e reparo sobre alertas do Infer em
Java. No reparo obtêm taxa lógica, isto é, correção equivalente à escrita por
um humano, de 97,30\% para desreferência nula e 91,77\% para vazamento de
recursos, sem ajuste fino. O mesmo grupo assina \citeonline{chatgpt-static}, e
os dois trabalhos compartilham dados, de modo que devem ser lidos como um só
conjunto de evidências.

\citeonline{sastgenius} integram o analisador Semgrep a um modelo Llama~3
ajustado, que faz triagem, geração de \textit{exploits} para validar alertas e
sugestão de correções. Trata-se de um \textit{preprint}, que relata a precisão
subindo de 35,7\% para 89,5\% mas não relata recall. Sua contribuição mais
relevante para o projeto apresentado aqui é apontar o risco de
\textit{prompt injection}, já que o código analisado, comentários inclusive,
entra no \textit{prompt} e pode tentar instruir o modelo a ignorar um defeito.

\subsection{O que a literatura de \llm\ já estabeleceu}

Três conclusões atravessam esses sete trabalhos. A divisão de trabalho que
funciona atribui ao analisador o raciocínio sobre o programa inteiro e ao
modelo o julgamento local, e nenhum dos trabalhos usa o modelo como analisador
autônomo. A política de decisão tem efeito maior que a escolha do modelo, pois os
filtros que descartam candidatos perdem de 6 a 25 pontos de recall
\cite{reducing-fp,adataint}, enquanto a política conservadora do LLift atinge
recall de 1,00 no mesmo ponto de integração \cite{llift}. E a arquitetura do
\textit{prompt} importa mais que o modelo escolhido, como mostra a ablação de
\citeonline{llift}, em que o mesmo modelo, sobre os mesmos dados, vai de 0,15 a
1,00 de recall ao ganhar ensino de regras do domínio por exemplos, pedidos
progressivos de contexto, decomposição da tarefa em conversas separadas e
autovalidação contra regras enviesadas para o lado seguro.

Nenhum desses sete trabalhos toca em concorrência, e nenhuma das quatro
ferramentas da Seção~\ref{sec:ferramentas} usa modelos de linguagem. A lacuna
que este trabalho ocupa é a interseção vazia entre as duas literaturas, e as
conclusões acima são o que se leva de uma para a outra.
